\runninghead{Scriptural Date and Location}
\section{Scriptural Date and Location}
\label{sec:date-location}

{\small Each passage the Mass reads or sings from Scripture, in canonical
order, with what is dated, where it belongs, and its date. Where several texts
of the Mass share one date, the row names them all. The Introit antiphon, the
three orations of the feast and the three of the Sunday have no biblical date.
The notes beneath give each passage's setting, attribution and dating
evidence.\par}
\medskip

\begin{chronologyorientationtable}
\chronologyorientationrow{row-e88e943a6faaa20b}{Two psalms of praise and of the fear of the Lord. The Clementine prints ``A psalm for David'' inside Ps 32:1, where the Hebrew has no title; Ps 33 is titled ``For David'' (33:1).}
\addlinespace
\chronologyorientationrow{row-bbfc04633a7fe6e8}{A hymn bidding the just rejoice and praise the Lord with the harp and a new canticle.}
\addlinespace
\chronologyorientationrow{row-fb6656afd6f4b4f0}{The Psalter, the book in which both psalms stand.}
\addlinespace
\chronologyorientationrow{row-d7a171509b4d3dc5}{David at the court of Achis, king of Geth, among the Philistines (1~Kings 21:10--15), the occasion the psalm's title names under the name Achimelech (Ps 33:1).}
\midrule
\chronologyorientationrow{row-43dce930fad06090}{The Book of Wisdom, whose answer to the wicked (Wis 2--3) the Offertory adapts; the proposed dates are reigns of Ptolemaic kings of Egypt.}
\midrule
\chronologyorientationrow{row-c1e1f88c6b954af2}{``A mountain'' (Mt 5:1), not named, after Jesus has gone ``about all Galilee'' and the multitudes have followed him (4:23--25).}
\addlinespace
\chronologyorientationrow{row-8b85083df79d6e95}{The Gospel according to St Matthew, from which the Alleluia verse, the Gospel and the Communion are taken.}
\midrule
\chronologyorientationrow{row-ca811ebc12efaaee}{John writes the visions he saw ``in the island which is called Patmos'' (Apoc 1:9).}
\end{chronologyorientationtable}

\subsection{Psalms 32 (33) and 33 (34): the Introit verse and the Gradual}
\label{sec:date-psalms}

The two psalms share the first and third rows of the table. Neither is dated
by an occasion of its own. Psalm~32 has no title in the Hebrew; the
Clementine Vulgate and the Douay print ``A psalm for David'' inside its first
verse, and the Missal's verse begins after it. Psalm~33 is titled ``For
David, when he changed his countenance before Achimelech, who dismissed him,
and he went his way''; the title answers to 1~Kings 21:10--15, where David,
fleeing Saul, comes to Achis, king of Geth, and feigns madness before him.
The attribution to David is the titles' own; the table marks it disputed and gives David's reign in the usual chronology from
the \work{Catholic Encyclopedia}'s article ``David'' (Corbett, vol.~4). The
reign is the reference point for that attribution, not a date of the psalms'
composition.

For Psalm~32 the chronology record has no historical setting, and the second
row says so. For Psalm~33 the table dates the episode the title names. The
figure is the year of the world that Haydock's Douay commentary gives at
1~Kings 21, reporting Ussher's reckoning; it is Ussher's date for
the flight to Achis, not an independent dating of the psalm. The title calls
the king Achimelech, while 1~Kings 21 calls him Achis; the episode is the
same, and the table follows the book of Kings for the place.

The composition date is the bound that the introduction to the Psalms in the
\work{New American Bible, Revised Edition} sets for the whole Psalter, in the
same sentence in which it holds that no individual psalm can be dated
securely. It dates the book, not these psalms.

\subsection{Wisdom 3:1--3: the Offertory}
\label{sec:date-wisdom}

The Book of Wisdom narrates no dated event in these verses; chapter~3 answers
the speech of the wicked in chapter~2. The two composition dates are those
that Francis Gigot's ``Book of Wisdom'' in the \work{Catholic Encyclopedia}
(vol.~15) records for the dispute over the book's date: in or after the time
of Ptolemy IV Philopator, and in or after the time of Ptolemy VII Physicon.
Both are marked disputed, and
the table gives neither as preferred. The dates are dates of the book, and
the antiphon is an adaptation of its verses, not the verses as the Vulgate
prints them.

\subsection{Matthew 5:1--12a, 5:8--10 and 11:28: the Gospel, the Communion and the Alleluia}
\label{sec:date-matthew}

The Sermon on the Mount is not dated in the chronology record, and the row
says so. Matthew places it on ``a mountain'' that he does not name, after
Jesus has gone about all Galilee and the multitudes have followed him from
Galilee, Decapolis, Jerusalem, Judea and beyond the Jordan (4:23--25). The
Alleluia verse and the Communion are drawn from the same Gospel and share its
composition row. The six composition dates are those of Eugène Jacquier's
``Gospel of St.~Matthew'' (vol.~10) and Alfred Durand's ``The New Testament''
(vol.~14) in the \work{Catholic Encyclopedia}; all are marked disputed, and the
table keeps the range from the early to the later dates rather than choosing
one. They are dates of the book, not of the sermon.

\subsection{Apocalypse 7:2--12: the Lesson}
\label{sec:date-apocalypse}

The two composition dates are those of Christian Van den Biesen's
``Apocalypse'' (vol.~1), which places the book in the reign of Domitian, and
of Alfred Durand's ``The New Testament'' (vol.~14), which gives an approximate
decade; both are marked disputed. John writes
that he was ``in the island which is called Patmos'' when he saw the visions
(1:9). The sealing and the multitude of chapter~7 are a vision and have no date
of their own.
